\chapter{Conclusion and Future Work}
\label{ch:conclusion}

\section{Conclusion}
This thesis designed, implemented, and empirically validated \textbf{Aetron}, a cross-platform Flutter mobile fitness application for tracking indoor and outdoor activities strictly with commodity smartphone sensors. Extending Aetron's pre-thesis prototype, we demonstrated that a transparent academic codebase can unify real-time geodetic state estimation, cross-platform environmental awareness, zero-velocity-hardened pedometry, crash-resilient local storage, and deterministic server auditing without relying on opaque proprietary black boxes.

Across five core objectives (\textbf{O1}--\textbf{O5}), our empirical findings establish concrete engineering and scientific advances:
\begin{enumerate}
    \item \textbf{Outdoor GNSS State Estimation \& Honest Negative Results:} Upgrading heuristic Fixed-Alpha blending ($6.722$\,m RMSE) to a Joseph-stabilized 4-state kinematic EKF with ZUPT ($3.560$\,m) and integrated \texttt{P01+P02} preprocessing ($3.432$\,m RMSE, \textbf{1.34\%} contiguous distance error on TU~Wien phase-RTK walks) eliminates stationary marker freeze and satisfies Scott et al.'s $<5\%$ \textbf{Good} benchmark---while matching Strava ($0.52$\,km) and scoring $2.5\%$ MAPE on an $800$\,m urban route. Preserving both the isolated AEKF \textit{NO-GO} ($+0.21\%$, CI crossing zero) and SnapperGPS~2 CEP50 cycling divergence ($138.226$\,m) cleanly documents sensor contract boundaries.
    \item \textbf{CAP-ALSTM Cross-Platform IOD \& Safety Isolation:} Implementing Bakirtzis et al.'s CAP-ALSTM~\cite{bakirtzis2022ciod} with Exact Shapley attribution and causal HMM filtering raises held-out Cambridge CIOD performance to \textbf{94.07\% Acc / 93.87\% Macro-F1} (7-channel Android without microphone), \textbf{78.82\% / 78.63\%} (4-channel mobile, $+9.04$\,pp Macro-F1 over Model~C-Mobile), and \textbf{76.14\% / 75.54\%} (3-channel iOS light-free, outperforming our ad-hoc iOS MLP+HMM baseline) with $\le 10^{-5}$ Python--Dart parity across $2{,}000$ windows. Transparent external evaluation on HAR-PMD ($35.29\%$--$42.47\%$ Macro-F1 without RSRP due to cross-device sensor-scale shift) alongside modest doorway F1 ($4.39\%$) confirms that \texttt{isObservationalOnly = true} remains the scientifically responsible default.
    \item \textbf{MV-GLRT Step Gating, Cadence, Stride, \& 2024 Compendium:} Combining adapted Skog MV-GLRT ZUPT and \texttt{StepMotionGate} removes $100\%$ of stationary hand-sway and screen-tap false steps while preserving $98.69\%$ running accuracy on Xu Group~3 and $93.16\%$ walking accuracy on Oxford WPD. Across $119{,}572$ video-verified steps, a $6$-second cadence window reduces flutter by $50.0\%$ ($8.41$\,spm MAE), Renaudin Eq.~8 LOSO improves walking/running stride error ($8.61$\,cm / $17.80$\,cm), offline sequential dwell clustering cuts false pauses $9.2\times$ ($98.97\%$ stop precision), Polyline saves $94.18\%$ bandwidth, and contiguous monotonic 2024 Compendium MET tables pass client--server parity and monotonicity regression gates across CI.
\end{enumerate}

\section{Limitations}
In keeping with transparent scientific reporting, six limitations remain:
\begin{enumerate}
    \item \textbf{Multi-Sport High-Speed RTK Gap:} Pedestrian walking was verified against Leica GS18 phase-RTK and cycling stressed on SnapperGPS~2 (lacking Doppler and native $1\sigma$ accuracy), but high-speed running and cycling still lack public native-smartphone RTK ground truth.
    \item \textbf{Co-Located Test Stand Trials:} TU~Wien Pixel~9a and Pixel~7~Pro runs share five physical routes on a collinear rig, representing paired environmental trials rather than 10 independent outdoor corridors.
    \item \textbf{Offline CIOD Coverage of iOS Barometer \& CoreLocation Extensions:} While the core \texttt{[MagIO+Acc+Daytime]} MLP + HMM is verified on $201{,}731$ CIOD epochs, fused \texttt{CMAltimeter} ($\Delta P_{10\text{s}}$) and \texttt{CoreLocation} blackout behavior requires dedicated live iPhone doorway logging.
    \item \textbf{Population-Average Calorimetry:} Without wearable heart-rate monitors or $\dot{V}\text{O}_2$ metabolic-cart calibration, 2024 Adult Compendium kcal remains a healthy-adult estimate.
    \item \textbf{Analytical vs.\ Physical Hardware Power Profiling:} CPU latency is directly benchmarked ($12.4$\,$\mu$s EKF; $0.78$\,ms CNN), whereas battery impact ($<0.05\%$/h) is modeled analytically rather than measured via hardware power monitors.
    \item \textbf{Emulator/Unit Tests vs.\ Endurance OEM Survival \& OSRM Absence:} Passing automated tests does not substitute for multi-hour survival under aggressive vendor task killers or formal SUS user evaluation, and route map matching reports unavailable until an OSRM server is connected.
\end{enumerate}

\section{Future Work}
The modular codebase and benchmark harness enable five immediate extensions:
\begin{enumerate}
    \item \textbf{Live iPhone Doorway \& Gym Ground-Truth Collection:} Capture synchronized iOS magnetometer, accelerometer, \texttt{CoreLocation}, barometer, and doorway ground-truth logs via \texttt{IodExperimentStorage} to validate graduation beyond observational mode.
    \item \textbf{Co-Mounted Geodetic RTK Rig for Running and Cycling:} Pair dual-frequency ZED-F9P RTK receivers with native iPhones and Android handsets across high-speed tracks.
    \item \textbf{Live OSRM HMM Road-Network Integration:} Attach a self-hosted OpenStreetMap OSRM graph behind \texttt{route-correction-worker} to activate production HMM map snapping~\cite{newson2009hidden}.
    \item \textbf{Individual Stride Calibration, Running Power, \& Wearable Bridge:} Support quick $100$\,m personal stride calibration, derive mechanical running power from elevation deltas, and optionally ingest Apple HealthKit / Health Connect heart-rate streams.
    \item \textbf{Formal Usability Study \& Endurance Battery Benchmarking:} Conduct a System Usability Scale (SUS) study alongside physical power-monitor profiling during long background sessions.
\end{enumerate}
